\documentclass{article}

\usepackage[a4paper, left=1.5cm, right=1.5cm, top=2cm, bottom=2cm]{geometry}

\usepackage{../../../../components/components}

\usepackage{fancyhdr}


% Configuration des en-têtes et pieds de page
\pagestyle{fancy}
\fancyhf{} % reset tout

\fancyhead[L]{DL2 Math-Info PR4}
\fancyhead[C]{Probabilités}
\fancyhead[R]{2025-2026}

\fancyfoot[L]{Ewen Rodrigues de Oliveira}
\fancyfoot[R]{\thepage}

\begin{document}

\docTitle{Chapitre 5 : Espérance}

\section{Définition et exemples}

\reminder{Rappel sur les séries. Soit $(a_n)_{n\in\mathbb{N}}$ une suite de nombres réels.\\
Si les $a_n$ sont positifs, alors :
\begin{itemize}
    \item la série $\sum_{n \geq 1} a_n$ est définie positive
    \item l'ordre dans lequel on somme les termes n'affecte pas la somme ni la finitude
    \item Si on partitionnel $\mathbb{N}^*$ en une collection dénombrable $(I_p)_{p\in\mathbb{N}}$ d'ensembles infinis, alors $\sum_{n \geq 1} a_n = \sum_{p \geq 1} \sum_{n \in I_p} a_n$
\end{itemize}
}

On suppose à présent que les $a_n$ sont de signes quelconques. 
Si $\sum_{n \ge 1} |a_n| < +\infty$ (série absolument convergente), alors :

\begin{enumerate}
    \item la série $\sum_{n \ge 1} a_n$ converge (dans $]-\infty, +\infty[$) ;
    
    \item l'ordre dans lequel l'on somme les $a_n$ n'influe pas sur la finitude ni sur la valeur de la série ;
    
    \item si l'on partitionne $\mathbb{N}^*$ en une collection dénombrable $(I_p)_{p \ge 1}$ de parties, alors
    \[
    \sum_{n \ge 1} a_n 
    = 
    \sum_{p \ge 1} \; \sum_{n \in I_p} a_n .
    \]
\end{enumerate}

\definition{
    Soit $X$ une variable aléatoire définie sur un espace probabilisé ($\Omega$, $\mathbb{P}$) et à valeurs dans un $\mathbb{R}_+$.
    L'\textbf{espérance de X} est définie comme la quantité :
\[\mathbb{E}[X] = \sum_{\omega \in \Omega} X(\omega) \cdot \mathbb{P}(\{\omega\}) \in [0, +\infty]\]
}

\definition{
    Soit $X : \Omega \to \mathbb{R}$ une variable aléatoire. On dit que $X$ est \textbf{intégrable} si $\mathbb{E}[|X|] < +\infty$.
}

\definition{
    Soit $X : \Omega \to \mathbb{R}$ une variable aléatoire intégrable.\\
    Posons $X^+ = \max(X, 0)$ et $X^- = \max(-X, 0)$. On définit l'\textbf{espérance de $X$} comme la quantité : $\mathbb{E}[X] = \mathbb{E}[X^+] - \mathbb{E}[X^-] \in \mathbb{R}$.
}

\theorem{Proposition}{Intégrabilité et espérance}{false}{
    Soit $X$ une variable aléatoire réelle intégrable.\\
    Alors $X$ admet une espérance et $\mathbb{E}[X] = \sum_{\omega \in \Omega} X(\omega) \cdot \mathbb{P}(\{\omega\})$.
}

\[
\begin{array}{c|c|c}
    & \text{Admet une espérance} & \text{Intégrable} \\
    \hline
    X \text{ positive et } \mathbb{E}[X] < +\infty & \text{oui (finie)} & \text{oui} \\
    X \text{ positive et } \mathbb{E}[X] = +\infty & \text{oui (infinie)} & \text{non} \\
    X \text{ quelconque et } \mathbb{E}[|X|] < +\infty & \text{oui (finie)} & \text{oui} \\
    X \text{ quelconque et } \mathbb{E}[|X|] = +\infty & \text{non} & \text{non}
\end{array}
\]
\remark{Dans le cas particulier où l'ensemble $E$ est fini, $X$ admet toujours une espérance et $X$ est même toujours intégrable.}

\example{
    On pose $\Omega = \{1, \ldots, 6\}$, $\mathbb{P}$ la probabilité uniforme et $X(\omega) = \omega$ la variable qui associe à chaque lancer de dé le nombre sur la face supérieure. $X$ prend un nombre fini de valeurs, donc $X$ est intégrable, et est à valeurs positives, donc $X$ admet une espérance. On trouve $\mathbb{E}[X] = \frac{1}{6} \cdot (1 + 2 + 3 + 4 + 5 + 6) = \frac{7}{2}$.
}

\vocabulary{En vocabulaire sur les séries, on dit que $X$ admet une espérance si la série $\sum_{x \in E} x \cdot \mu_X(\{x\})$ est absolument convergente.}

\section{Formule de transfert et propriétés}

\subsection{Énoncés}

\theorem{Théorème}{Formule de transfert}{false}{
    Soit $X$ une variable aléatoire sur un espace probabilisé $(\Omega, \mathbb{P})$ à valeurs dans un ensemble dénombrable $E$.\\
    Soit $g \colon E \to F$ une fonction, avec $F \subset \mathbb{R}$ un ensemble dénombrable.
    \begin{enumerate}
        \item Si $g$ est à valeurs positives, alors :
        \[ \mathbb{E}[g(X)] = \sum_{x \in E} g(x) \cdot \mu_X(\{x\}) .\]
        \item Si $g$ est de signe quelconque, alors : la variable aléatoire $g(X)$ est intégrable si et seulement si la série $\sum_{x \in E} |g(x)| \cdot \mu_X(\{x\})$ est finie, et dans ce cas :
        \[ \mathbb{E}[g(X)] = \sum_{x \in E} g(x) \cdot \mu_X(\{x\}) .\]
    \end{enumerate}
}

\theorem{Corollaire}{Expression de l'espérance à partir de la loi de $X$}{false}{
    Soit $X$ une variable aléatoire réelle discrète dans un ensemble $E$.
    \begin{enumerate}
        \item Si $X$ est à valeurs positives, alors :
        \[ \mathbb{E}[X] = \sum_{\omega \in \Omega} X(\omega) \cdot \mathbb{P}(\{\omega\}) .\]
        \item Si $X$ est de signe quelconque, alors : $X$ est intégrable si et seulement si la série $\sum_{\omega \in \Omega} |X(\omega)| \cdot \mathbb{P}(\{\omega\})$ est finie, et dans ce cas :
        \[ \mathbb{E}[X] = \sum_{\omega \in \Omega} X(\omega) \cdot \mathbb{P}(\{\omega\}) .\]
    \end{enumerate}
}

\subsection{Conséquences}

On a donc vu que l'espérance d'une variable aléatoire discrète est une série de la forme $\sum_{n \ge 1} a_n$ avec $a_n = x \cdot \mu_X(\{x\})$ ou $a_n = X(\omega) \cdot \mathbb{P}(\{\omega\})$ qui converge absolument. On peut alors appliquer les propriétés des séries absolument convergentes pour obtenir les propriétés suivantes de l'espérance.
\theorem{Propriété}{}{false}{
    \begin{enumerate}
        \item Linéarité : pour toutes variables aléatoires intégrables $X, Y$ et tous réels $a, b$, la variable aléatoire $aX + bY$ est intégrable et $\mathbb{E}[aX + bY] = a \mathbb{E}[X] + b \mathbb{E}[Y]$.
        \item Monotonie : pour toutes variables aléatoires $X, Y$ admettant une espérance, si $X(\omega) \leq Y(\omega)$ pour tout $\omega \in \Omega$, alors $\mathbb{E}[X] \leq \mathbb{E}[Y]$.
        \item Valeur absolue : $|\mathbb{E}[X]| \leq \mathbb{E}[|X|]$ pour toute variable aléatoire $X$ admettant une espérance.
    \end{enumerate}
}

\theorem{Proposition}{}{false}{
    Soit $X$ une variable aléatoire discrète. Supposons $X(\omega) \geq 0$ pour tout $\omega \in \Omega$ (ie. $X$ est à valeurs positives).\\
    Alors $E[X] = 0 \Leftrightarrow \mathbb{P}(X = 0) = 1$.\\\\

    Attention, si $X$ est de signe quelconque, $\Rightarrow$ n'est plus vérifié.
}

\theorem{Proposition}{Variable aléatoire entière}{false}{
    Soit $X$ une variable aléatoire à valeurs dans $\mathbb{N}$. Alors :
    \[ \mathbb{E}[X] = \sum_{k=1}^{+\infty} \mathbb{P}(X \geq k) .\]
}

\example{
    Soit $X$ une variable aléatoire qui suit le nombre de la face supérieure d'un lancer de dé. On trouve $\mathbb{E}[X] = \sum_{k=1}^{+\infty} \mathbb{P}(X \geq k) = \sum_{k=1}^{6} \frac{7-k}{6} = \frac{7}{2}$.
}

\theorem{Proposition}{Inégalité de Markov}{false}{
    Si $X \geq 0$ et $a > 0$, alors :
    \[ \mathbb{P}(X \geq a) \leq \frac{\mathbb{E}[X]}{a} .\]
}

\example{
    Soit $X$ une variable aléatoire qui suit le nombre de la face supérieure d'un lancer de dé. On trouve $\mathbb{P}(X \geq 5) = \frac{1}{3} \leq \frac{\mathbb{E}[X]}{5} = \frac{7}{10}$.
}

\remark{Cette inégalité est très utile pour majorer la probabilité que $X$ prenne des valeurs grandes à partir de l'espérance de $X$. Elle est utilisée notamment dans la preuve de l'inégalité de Bienaymé-Tchebychev, qui permet de majorer la probabilité que $X$ prenne des valeurs éloignées de son espérance à partir de la variance de $X$.}

\section{Variance}
\subsection{Définition et propriétés}
\vocabulary{
    Soit $X$ une variable aléatoire réelle discrète. On dit que $X$ est de \textbf{carré intégrable} si $X^2$ est intégrable, \textit{i.e.} si $\mathbb{E}[X^2] < +\infty$.
}

\theorem{Lemme}{}{false}{
    Soit $X$ une variable aléatoire réelle discrète de carré intégrable. Alors $X$ est intégrable.\\
    Si de plus, $Y$ est une variable aléatoire réelle discrète de carré intégrable, alors $X + Y$ est de carré intégrable.
}

\definition{
    Soit $X$ une variable aléatoire réelle discrète de carré intégrable. On définit la variance de $X$ comme la quantité :
    \[ \mathrm{Var}(X) = \mathbb{E}[(X - \mathbb{E}[X])^2] .\]
}

\remark{La variance de $X$ est la moyenne des carrés des écarts à l'espérance de $X$. Elle mesure la dispersion de $X$ autour de son espérance : plus la variance est grande, plus les valeurs de $X$ sont éloignées de son espérance.}

\theorem{Proposition}{}{false}{
    Soit $X$ une variable aléatoire réelle discrète de carré intégrable. Alors :
    \[ \mathrm{Var}(X) = \mathbb{E}[X^2] - (\mathbb{E}[X])^2 \; \in [0, +\infty[ .\]
}

\theorem{Corollaire}{}{false}{
    Soit $X$ une variable aléatoire réelle discrète de carré intégrable. Alors $\mathrm{Var}(X) = 0 \Leftrightarrow \mathbb{P}(X = \mathbb{E}[X]) = 1$, \textit{i.e. $X$ est constante}.
}

\remark{Si $X$ suite une unité précise, alors on a un problème : l'unité de la variance est le carré de l'unité de $X$. Par exemple, si $X$ est mesuré en mètres, alors $\mathrm{Var}(X)$ est mesurée en mètres carrés. D'où la définition qui suit.}

\definition{
    Soit $X$ une variable aléatoire réelle discrète de carré intégrable. On définit l'écart-type de $X$ comme la quantité :
    \[ \sigma(X) = \sqrt{\mathrm{Var}(X)} .\]
}

\subsection{Application aux lois usuelles}

Pour les lois usuelles vu au chapitre 4 (Variables aléatoires), on trouve les résultats suivants :

\begin{center}
\renewcommand{\arraystretch}{2.5}
\begin{tabular}{|l|c|c|c|}
\hline
\textbf{Loi / Paramètres} & \textbf{$\mathbb{P}(X=k)$} & \textbf{Espérance} & \textbf{Variance} \\
\hline
Uniforme $\mathcal{U}(\{1, \ldots, n\})$ sur $\{1,\ldots,n\}$ & $\dfrac{1}{n}$ & $\dfrac{n+1}{2}$ & $\dfrac{n^2-1}{12}$ \\
\hline
Bernoulli $Ber(p)$ sur $\{0,1\}$ & $p^k(1-p)^{1-k}$ & $p$ & $p(1-p)$ \\
\hline
Binomiale $\mathcal{B}(n,p)$ sur $\{0,\ldots,n\}$ & $\binom{n}{k}p^k(1-p)^{n-k}$ & $np$ & $np(1-p)$ \\
\hline
Poisson $\mathcal{P}(\lambda)$ sur $\mathbb{N}$ & $e^{-\lambda}\dfrac{\lambda^k}{k!}$ & $\lambda$ & $\lambda$ \\
\hline
Géométrique $\mathcal{G}(p)$ sur $\mathbb{N}$ & $p(1-p)^{k}$ & $\dfrac{1}{p}$ & $\dfrac{1-p}{p^2}$ \\
\hline
Hypergéométrique $\mathcal{H}$ sur $\mathcal{A}$ & $\dfrac{\binom{K}{k}\binom{N-K}{n-k}}{\binom{N}{n}}$ & $\dfrac{nK}{N}$ & $n\dfrac{K}{N}\left(1-\dfrac{K}{N}\right)\dfrac{N-n}{N-1}$ \\
\hline
\end{tabular}
\renewcommand{\arraystretch}{1}
\end{center}

\smallskip
\noindent
Ici, $\mathcal{A} = \{\max(0,n-(N-K)),\ldots,\min(n,K)\} = \{k \in \mathbb{N} : k \leq K, n-k \leq N-K\}$.

\section{Espérance et indépendance}

Il existe des liens entre l'espérance et l'indépendance de variables aléatoires.

\theorem{Proposition}{Espérance du produit de variables aléatoires indépendantes}{false}{
    Soient $X$ et $Y$ deux variables aléatoires à valeurs dans $E$ et $F$ deux ensembles dénombrables respectivement.
    \begin{enumerate}
        \item Si $E,F \subset \mathbb{R}$ et si $X$ et $Y$ sont intégrables, alors $XY$ est intégrable et on a : \[ \mathbb{E}[XY] = \mathbb{E}[X] \cdot \mathbb{E}[Y] .\]
        \item \textit{Plus généralement}, si $f \colon E \to \mathbb{R}$ et $g \colon F \to \mathbb{R}$ sont des fonctions telles que $f(X)$ et $g(Y)$ sont intégrables, alors $f(X)g(Y)$ est intégrable et on a : \[ \mathbb{E}[f(X)g(Y)] = \mathbb{E}[f(X)] \cdot \mathbb{E}[g(Y)] .\]
    \end{enumerate}
}

\theorem{Proposition}{Espérance du produit de $N$ variables aléatoires indépendantes}{false}{
    Soit $N \geq 1$ un entier et soient $(X_k)_{1 \leq k \leq N}$ une collection de variables aléatoires indépendantes.\\
    On suppose que chaque variable $X_k$ est à variable dans ensemble dénombrable $E_k$.\\\\

    $\forall f_1, \ldots, f_N : E_1 \to \mathbb{R}, \ldots, E_N \to \mathbb{R}$ telles que $f_k(X_k)$ est intégrable pour tout $k \leq N$, alors $\prod_{k=1}^N f_k(X_k)$ est intégrable et :
    \[ \mathbb{E}\left[ \prod_{k=1}^N f_k(X_k) \right] = \prod_{k=1}^N \mathbb{E}[f_k(X_k)] .\]
}

\remark{Ici, on tire les conséquences de l'indépendance de $X$ et $Y$ pour obtenir une égalité sur les espérances. On peut aussi faire le chemin inverse : à partir d'une égalité sur les espérances, on peut tirer des conséquences sur l'indépendance de $X$ et $Y$. C'est ce que fait la proposition suivante.}

\theorem{Proposition}{Caractérisation de l'indépendance}{false}{
    Soit $N \geq 1$ et soit $(X_k)_{1 \leq k \leq N}$ une collection de variables aléatoires réelles discrètes.\\
    On suppose que chaque variable $X_k$ est à variable dans ensemble dénombrable $E_k$.\\\\

    Ces variables sont indépendantes si et seulement si pour toutes fonctions $f_1, \ldots, f_N$ \textbf{bornées} de $E_1, \ldots, E_N$ dans $\mathbb{R}$, on a :
    \[ \mathbb{E}\left[ \prod_{k=1}^N f_k(X_k) \right] = \prod_{k=1}^N \mathbb{E}[f_k(X_k)] .\]
}


\newpage
\tableofcontents

\end{document}